\cleardoublepage
\chapter{SKEMA \textit{KNOWLEDGE GRAPH} DAN CONTOH DATA}
\label{lampiran:skema-data}

Lampiran ini menyajikan bentuk konkret data yang mengalir dalam \textit{pipeline ingestion} sistem \textit{GraphRAG} Kubernetes: dari definisi \texttt{swagger.json} sebagai masukan, menjadi \textit{node} dan \textit{relationship} di \textit{knowledge graph} Neo4j sebagai keluaran, sebagaimana dibahas pada Bab~\ref{chap:implementasi}. Bagian terakhir merangkum skema graf secara menyeluruh: label \textit{node} dan seluruh 18 tipe \textit{edge} yang digunakan sistem.

% ─────────────────────────────────────────────────────────────────────────────
\section{Contoh Masukan: Cuplikan \texttt{definitions.json}}
% ─────────────────────────────────────────────────────────────────────────────

Cuplikan berikut menunjukkan dua entri OpenAPI \textit{definitions} dari \texttt{swagger.json} Kubernetes sebagaimana diproses oleh \textit{SwaggerGraphBuilder} pada Fase 1 (Bab~\ref{chap:implementasi}): \texttt{DeploymentStrategy} dan \texttt{RollingUpdateDeployment}, yang saling terhubung melalui \textit{field} \texttt{rollingUpdate}.

\lstinputlisting[
    language=json,
    frame=single,
    captionpos=t,
    caption={\textit{Cuplikan definitions.json} — \texttt{Deployment\allowbreak{}Strategy} dan \texttt{Rolling\allowbreak{}Update\allowbreak{}Deployment}},
    label={lst:kg-input-swagger}
]{listings/kg-input-swagger-snippet.json}

% ─────────────────────────────────────────────────────────────────────────────
\section{Contoh Keluaran: Node dan Relationship di Neo4j}
% ─────────────────────────────────────────────────────────────────────────────

Cuplikan di atas, setelah melalui Fase 1 (pembuatan \textit{node}) dan Fase 2 (pembuatan \textit{edge} struktural), menghasilkan dua \textit{node} berlabel \texttt{Definition} beserta satu \textit{relationship} \texttt{HAS\_PROPERTY} yang menghubungkannya. Properti \texttt{embedding} (vektor 1536 dimensi dari \texttt{text-embedding-3-small}, Fase 1.5) dipangkas untuk keterbacaan.

\lstinputlisting[
    language=json,
    frame=single,
    captionpos=t,
    caption={Representasi \textit{Node} dan \textit{Relationship} Hasil \textit{Ingestion} di Neo4j},
    label={lst:kg-output-node}
]{listings/kg-output-node-example.json}

% ─────────────────────────────────────────────────────────────────────────────
\section{Skema Graf: Label Node dan Tipe Edge}
% ─────────────────────────────────────────────────────────────────────────────

\begin{sloppypar}
Seluruh \textit{node} di \textit{knowledge graph} berlabel \texttt{Definition}; \textit{node} yang merepresentasikan \textit{resource} Kubernetes tingkat atas (memiliki \textit{Group-\allowbreak{}Version-\allowbreak{}Kind} pada \texttt{x-kubernetes-group-\allowbreak version-kind}) mendapat label tambahan \texttt{K8sResource}. Tabel~\ref{tbl:edge-schema} merangkum seluruh 18 tipe \textit{edge} yang dibangun pada Fase 2 (struktural), Fase 2.5 (\textit{inheritance}/\textit{union type}), dan Fase 3 (semantik) \textit{pipeline ingestion}.
\end{sloppypar}

\begin{footnotesize}
\begin{longtable}{|p{3.4cm}|p{4.4cm}|p{4.5cm}|}
\caption{Skema 18 Tipe \textit{Edge} \textit{Knowledge Graph}}
\label{tbl:edge-schema} \\
\hline
\textbf{Tipe Edge} & \textbf{Contoh Relasi} & \textbf{Keterangan} \\
\hline
\endfirsthead
\hline
\multicolumn{3}{l}{Tabel \ref{tbl:edge-schema} Skema 18 Tipe \textit{Edge} \textit{Knowledge Graph} (lanjutan)} \\
\hline
\textbf{Tipe Edge} & \textbf{Contoh Relasi} & \textbf{Keterangan} \\
\hline
\endhead
\hline
\multicolumn{3}{r}{\textit{bersambung ke halaman berikutnya}} \\
\endfoot
\hline
\endlastfoot
\textit{HAS\_PROPERTY} & Resource $\rightarrow$ \textit{field}/tipe anak & Relasi struktural utama; setiap \textit{field} objek pada skema OpenAPI \\
\hline
\textit{EXTENDS} & Deployment $\rightarrow$ DeploymentSpec & Pola pewarisan implisit \textit{resource} $\rightarrow$ \texttt{*Spec}-nya \\
\hline
\textit{ONE\_OF} & \texttt{IntOrString} $\rightarrow$ opsi tipe & Alternatif tipe polimorfik (\textit{union type}) \\
\hline
\textit{ANY\_OF} & \texttt{EnvVarSource} $\rightarrow$ opsi tipe & Alternatif tipe polimorfik (\textit{union type}) \\
\hline
\textit{SCALES\_RESOURCE} & HPA $\rightarrow$ Deployment\slash StatefulSet\slash ReplicaSet & \textit{Horizontal Pod Autoscaler} menyasar target skalanya \\
\hline
\textit{CONTAINS\_\allowbreak POD\_TEMPLATE} & Deployment\slash StatefulSet\slash DaemonSet\slash Job $\rightarrow$ PodTemplateSpec & \textit{Workload controller} berisi templat Pod \\
\hline
\textit{CONTAINS\_\allowbreak JOB\_TEMPLATE} & CronJob $\rightarrow$ JobTemplateSpec & CronJob berisi templat Job \\
\hline
\textit{BINDS\_ROLE} & RoleBinding\slash ClusterRoleBinding $\rightarrow$ Role\slash ClusterRole & Relasi RBAC \\
\hline
\textit{BINDS\_\allowbreak SERVICE\_ACCOUNT} & RoleBinding\slash ClusterRoleBinding $\rightarrow$ ServiceAccount & Relasi RBAC \\
\hline
\textit{HAS\_CONTAINER} & PodSpec $\rightarrow$ Container & PodSpec berisi satu atau lebih Container \\
\hline
\textit{CLAIMS\_VOLUME} & StatefulSet $\rightarrow$ PersistentVolumeClaim & \textit{Volume claim templates} StatefulSet \\
\hline
\textit{MOUNTS\_VOLUME} & PodSpec $\rightarrow$ Volume & Volume yang dipasang pada Pod \\
\hline
\textit{USES\_STORAGE\_CLASS} & PersistentVolumeClaim $\rightarrow$ StorageClass & PVC menyasar kelas penyimpanan \\
\hline
\textit{LOADS\_CONFIGMAP} & Container $\rightarrow$ ConfigMap & Konfigurasi dimuat via \texttt{envFrom}/\texttt{volumeMounts} \\
\hline
\textit{USES\_SECRET} & PodSpec $\rightarrow$ Secret & \textit{Secret} digunakan oleh Pod \\
\hline
\textit{SELECTS\_POD} & Service $\rightarrow$ Pod & Service memilih Pod melalui \textit{label selector} \\
\hline
\textit{ROUTES\_TO\_SERVICE} & Ingress $\rightarrow$ Service & Ingress merutekan \textit{traffic} ke Service \\
\hline
\textit{USES\_\allowbreak SERVICE\_ACCOUNT} & PodSpec $\rightarrow$ ServiceAccount & Pod berjalan dengan identitas ServiceAccount \\
\hline
\end{longtable}
\end{footnotesize}
